\scanpage{57}
\mspage{40}
\margin{n°382}
est formé des couples $(h,k)\in T_x(X)\times T_y(Y)$ tels que $T_x(f).h=T_y(g).k$.

Le morphisme $g':X\times_S Y\longrightarrow X$ est parfois appelé \uline{l'image réciproque} du morphisme $g:Y\longrightarrow S$ par f ; par abus de langage, on dit aussi que $X\times_S Y$ est \uline{l'image réciproque de} Y \uline{par} f. Si g est une submersion (resp. une immersion), il en est de même de g'. Si X est une sous-variété de S, et si f est l'injection canonique de X dans S, le couple $(f,g)$ est transversal si et seulement si g est transversal à X (cf. n° 4.11.1) ; de plus, f' induit un isomorphisme de $X\times_S Y$ sur la sous-variété $g^{-1}(X)$ de Y (ce qui justifie, dans ce cas, le terme d' ``image réciproque'' appliqué à $X\times_S Y$). En particulier, $S\times_S Y$ s'identifie à Y.

4.11.4. Soit $f:X\longrightarrow S$, $g:Y\longrightarrow S$ un couple transversal de morphismes. Soit T une variété, et soient $u:T\longrightarrow X$, $v:T\longrightarrow Y$ des morphismes tels que $f\circ u=g\circ v$. Il existe alors un morphisme $j:T\longrightarrow X\times_S Y$ et un seul tel que $u=g'\circ j$, $v=f'\circ j$, les notations étant celles du n° précédent. Le morphisme j se note $(f,g)_S$.

4.12 -- \uline{Variétés de groupes}.

4.12.1. Soit G un groupe. Une structure de variété sur G est dite \uline{compatible} avec la structure de groupe de G si l'application $(x,y)\longrightarrow xy$ de $G\times G$ dans G est un morphisme, ainsi que l'application $x\longrightarrow x^{-1}$ de G dans lui-même. Muni de ces deux structures, G est appelé une \uline{variété de groupe} (``de classe $C^r$'' si l'on veut préciser) ; pour $r=\omega$ (resp. $r\ne\omega$), on dit aussi ``groupe analytique'' (resp. ``groupe différentiable''). Toute variété de groupe est pure. On appelle \uline{morphisme de variétés de groupe} (ou simplement \uline{morphisme}) toute application d'une variété de groupe dans une autre qui est à la fois une \uline{représentation} de groupes et un \uline{morphisme} de variétés.

Si G est une variété de groupe, la structure topologique sous-jacente à la structure de variété de G en fait un \uline{groupe topologique métrisable et complet} ; il est localement compact si K est localement compact.
\mspage{56}
